\chapter{Server concorrenti: processi e thread}

Abbiamo un server che funziona e che non si può usare: serve un client, e tutti gli altri aspettano che quello se ne vada.
\begin{itemize}
    \item Nel frattempo abbiamo smontato i tre meccanismi che lo risolvono: un \textbf{secondo processo}, un \textbf{secondo
          thread}, e un \textbf{mutex} per rendere sicura la condivisione.
    \item Nessuno dei tre ha mai toccato un socket: la \texttt{fork} stampava una riga due volte, i thread incrementavano un
          contatore.
    \item Ora entrano nel ciclo di \texttt{accept}, e lo stesso servizio si scrive due volte: una con i processi, una con i
          thread.
    \item Le due versioni differiscono in una cosa sola, che poi decide tutto il resto: se il codice che serve un client
          \textbf{può vedere gli altri}.
\end{itemize}

\subsection{Che cosa rende concorrente un server}

\begin{itemize}
    \item Il kernel accetta già le connessioni mentre siamo occupati: restano nella \textbf{coda di listen}. Quello che non
          può fare è servirle.
    \item La coda quindi compra \textbf{spazio, non tempo}. Un client che ci sta dentro aspetta esattamente quanto prima, e
          quando la coda è piena il successivo viene rifiutato.
    \item E il server, quasi sempre, non sta calcolando: un server di chat passa la vita bloccato in \texttt{read}.
\end{itemize}

\dfn{Server concorrente}{
    Un server è \textbf{concorrente} quando il tempo che un client passa a pensare non diventa l'attesa di un altro: il
    secondo viene servito mentre il primo tace.
}

Ci sono tre modi di arrivarci: \textbf{un processo per client}, \textbf{un thread per client}, e \textbf{un solo flusso di
controllo che osserva molte connessioni}. In quest'ordine: i primi due in questo capitolo, il terzo nel prossimo.

% ══════════════════════════════════════════════════════════════
\section{Un processo per client}
% ══════════════════════════════════════════════════════════════

Nel meccanismo non c'è nulla di nuovo: dopo \texttt{accept} il server fa \texttt{fork}, e il figlio eredita
\textbf{entrambi} i socket, quello di ascolto e quello connesso. È la figura dei descriptor ereditati del capitolo sui
processi: ciascuno dei due chiude quello che non usa, e ogni oggetto vive finché non se ne va il suo ultimo descriptor.
Cambia solo ciò che va preparato \textbf{prima} del ciclo.

\begin{lstlisting}[language={[POSIX]C}, style=wnumbers, title={Il server concorrente: che cosa cambia prima del ciclo}]
static void reap(int sig) {                 /* installato una volta, prima di ogni fork */
    int saved = errno;
    while (waitpid(-1, NULL, WNOHANG) > 0)
        ;                                   /* tutti i figli che hanno finito */
    errno = saved;
}

int main(int argc, char *argv[]) {
    struct sigaction sa;
    int listen_fd = listening_socket(argc > 1 ? argv[1] : "5000");

    sa.sa_handler = reap;
    sigemptyset(&sa.sa_mask);
    sa.sa_flags = 0;
    sigaction(SIGCHLD, &sa, NULL);          /* altrimenti i figli si accumulano come zombie */
    signal(SIGPIPE, SIG_IGN);               /* un client che se ne va non deve ucciderci */
\end{lstlisting}

Due righe, ed entrambe riguardano \textbf{qualcun altro che finisce}: un figlio che è morto, e un client che ha chiuso
mentre gli stavamo scrivendo.

\begin{lstlisting}[language={[POSIX]C}, style=wnumbers, title={Il ciclo di accept che fa fork}]
    for (;;) {
        int fd = accept(listen_fd, NULL, NULL);
        if (fd < 0) {
            if (errno == EINTR) continue;   /* e' morto un figlio mentre aspettavamo */
            perror("accept"); continue;
        }
        pid_t pid = fork();
        if (pid < 0) { perror("fork"); close(fd); continue; }   /* si rifiuta solo questo */

        if (pid == 0) {                     /* il figlio: solo questa connessione */
            close(listen_fd);
            serve(fd);
            close(fd);
            _exit(0);                       /* non exit(): i buffer non sono suoi */
        }
        close(fd);                          /* il genitore torna subito in accept */
    }
}
\end{lstlisting}

\begin{itemize}
    \item L'\texttt{EINTR} è una conseguenza diretta dell'handler: il \texttt{SIGCHLD} arriva mentre il genitore è
          bloccato in \texttt{accept}, e la chiamata ritorna con $-1$. Non è un errore: si ricomincia.
    \item Una \texttt{fork} che fallisce (troppi processi, poca memoria) rifiuta \textbf{questo} client, non ferma il
          server.
\end{itemize}

\begin{lstlisting}[language={[POSIX]C}, style=wnumbers, title={Il figlio: il servizio non è cambiato}]
static unsigned long served = 0;            /* quanti messaggi abbiamo inoltrato */

static void serve(int fd) {                 /* un'intera conversazione */
    char buf[256], head[32];
    ssize_t n;
    while ((n = read(fd, buf, sizeof buf)) > 0) {
        int h = snprintf(head, sizeof head, "[%lu] ", ++served);
        if (write(fd, head, h) != h) break;
        if (write(fd, buf, n) != n) break;  /* quello che e' arrivato: non ancora una riga */
    }
    if (n < 0) perror("read");
}
\end{lstlisting}

La concorrenza sta \textbf{tutta nel ciclo}, non nel servizio: questa funzione è quella del server sequenziale, intatta.
Teniamo d'occhio \texttt{served}.

\subsection{Venti client insieme}

\begin{lstlisting}[language=bash, title={Venti client, venti figli}]
$ ./chatd 5000 &
$ for i in $(seq 1 20); do sleep 60 | nc 127.0.0.1 5000 & done
$ ps --ppid $(pgrep -o chatd) -o pid,stat,cmd | head -4
    PID STAT CMD
   7413 S    ./chatd 5000       <- un figlio per ogni client connesso
   7414 S    ./chatd 5000
   7415 S    ./chatd 5000
$ ss -Htn state established '( sport = :5000 )' | wc -l
20
\end{lstlisting}

\texttt{pgrep -o} sceglie il processo più vecchio con quel nome, cioè il genitore; \texttt{--ppid} elenca i suoi figli.
Venti client, venti figli, e ciascuno servito. Allora scriviamo una riga in uno di loro: chi la vede?

\warning{Nessuno, tranne chi l'ha mandata}{
    Il client 1 invia una riga, e altri diciannove sono connessi. Che cosa vedono? \textbf{Niente, mai}. La riga è stata letta
    da un figlio, e quel figlio non ha alcun descriptor verso le altre connessioni: il genitore le ha chiuse una per una, e i
    fratelli sono stati creati prima che la sua esistesse.
    \begin{itemize}
        \item Per lo stesso motivo \texttt{served} è inutile: ogni client si sente dire che il suo è il messaggio numero
              uno, perché ogni figlio ha ricevuto la propria copia al momento della \texttt{fork}.
        \item Abbiamo costruito venti server di eco indipendenti che per caso condividono una porta: va benissimo per un
              server web, è inutile per una chat.
    \end{itemize}
}

Gli spazi di indirizzi separati sono esattamente ciò che abbiamo chiesto. Se i processi devono parlarsi,
\textbf{deve portarlo il kernel}.

% ══════════════════════════════════════════════════════════════
\section{Far parlare i processi}
% ══════════════════════════════════════════════════════════════

\begin{figure}[H]
\centering
\begin{tikzpicture}[every node/.style={font=\scriptsize\sffamily},
    pr/.style={draw=netblue!70!black, fill=netblue!8, rounded corners=2pt, minimum width=1.9cm, minimum height=0.7cm},
    a/.style={-{Stealth[length=1.8mm]}, thick},
    aa/.style={{Stealth[length=1.8mm]}-{Stealth[length=1.8mm]}, thick}]
  \node[font=\small\bfseries\sffamily] at (0,1.5) {scambio di messaggi};
  \node[pr] (a1) at (-1.5,0.6) {processo A};
  \node[pr] (b1) at (1.5,0.6)  {processo B};
  \node[draw=black!50, fill=black!5, rounded corners=2pt, minimum width=5.0cm, minimum height=0.6cm] (k1) at (0,-0.9)
    {kernel: il buffer della pipe o del socket};
  \draw[a] (a1.south) -- node[left]{\texttt{write}} (a1.south |- k1.north);
  \draw[a] (b1.south |- k1.north) -- node[right]{\texttt{read}} (b1.south);
  \node[text=black!60] at (0,-1.6) {una chiamata di sistema per ogni scambio};
  \node[font=\small\bfseries\sffamily] at (7.2,1.5) {memoria condivisa};
  \node[pr] (a2) at (5.7,0.6) {processo A};
  \node[pr] (b2) at (8.7,0.6) {processo B};
  \node[draw=netgreen!70!black, fill=netgreen!10, rounded corners=2pt, minimum width=5.0cm, minimum height=0.6cm] (m) at
    (7.2,-0.9) {una pagina mappata da entrambi};
  \draw[aa] (a2.south) -- node[left]{load, store} (a2.south |- m.north);
  \draw[aa] (b2.south) -- node[right]{load, store} (b2.south |- m.north);
  \node[text=black!60] at (7.2,-1.6) {il kernel interviene una volta, poi coordinarsi tocca a noi};
\end{tikzpicture}
\caption{I due modi in cui due processi si parlano. A sinistra ogni dato passa per il kernel, che ne garantisce anche
         l'ordine; a destra il kernel prepara la pagina e poi si fa da parte.}
\end{figure}

\begin{itemize}
    \item \textbf{Scambio di messaggi}: ogni scambio è una chiamata di sistema, il kernel tiene i dati, e l'ordine di arrivo
          è una sua garanzia.
    \item \textbf{Memoria condivisa}: il kernel viene interpellato una volta sola, poi i due processi leggono e scrivono
          memoria ordinaria, e il coordinamento è \textbf{interamente loro}.
    \item La seconda è più veloce, e ci riporta a \texttt{count++}: la corsa critica è la stessa che hanno i thread.
\end{itemize}

\subsection{La pipe}

\dfn{Pipe}{
    \texttt{pipe(fd)} restituisce \textbf{due descriptor su un unico buffer del kernel}: \texttt{fd[0]} per leggere e
    \texttt{fd[1]} per scrivere. Dopo una \texttt{fork} entrambi i processi li hanno entrambi.
}

\begin{figure}[H]
\centering
\begin{tikzpicture}[every node/.style={font=\scriptsize\sffamily},
    fd/.style={draw=black!70, fill=white, minimum width=1.6cm, minimum height=0.45cm, font=\scriptsize\ttfamily},
    fdx/.style={fd, dashed, draw=netred, text=netred},
    a/.style={-{Stealth[length=1.8mm]}, thick},
    ax/.style={a, dashed, netred}]
  \node[font=\small\bfseries\sffamily, anchor=east] at (-4.2,1.3) {genitore};
  \node[font=\small\bfseries\sffamily, anchor=east] at (-4.2,-1.3) {figlio};
  \node[fd]  (g1) at (-3,1.3)  {fd[1]};
  \node[fdx] (g0) at (3,1.3)   {fd[0]};
  \node[fdx] (f1) at (-3,-1.3) {fd[1]};
  \node[fd]  (f0) at (3,-1.3)  {fd[0]};
  \node[draw=netblue!70!black, fill=netblue!8, rounded corners=2pt, minimum width=3.4cm, minimum height=0.7cm] (p) at (0,0)
    {buffer nel kernel};
  \draw[-{Stealth[length=2mm]}, netblue, thick] (-1.3,-0.22) -- (1.3,-0.22);
  \draw[a]  (g1.south) |- ([yshift=1.5mm]p.west);
  \draw[ax] (f1.north) |- ([yshift=-1.5mm]p.west);
  \draw[ax] ([yshift=1.5mm]p.east) -| (g0.south);
  \draw[a]  ([yshift=-1.5mm]p.east) -| (f0.north);
  \node[text=black!60, anchor=south] at (-3,1.6) {lato di scrittura};
  \node[text=black!60, anchor=south] at (3,1.6) {lato di lettura};
\end{tikzpicture}
\caption{Una pipe dopo la fork, quando il genitore scrive e il figlio legge. In rosso tratteggiato i due descriptor da
         chiudere: ciascun processo rinuncia al lato che non usa.}
\end{figure}

È \textbf{a senso unico} ed è \textbf{fra parenti}: l'unico modo di entrarci è ereditare un descriptor che esiste già.

\begin{lstlisting}[language={[POSIX]C}, style=wnumbers, title={Chiudere i lati che non si usano}]
int fd[2];
pipe(fd);                             /* controllo d'errore omesso */
if (fork() == 0) {
    close(fd[1]);                     /* chi legge rinuncia al lato di scrittura */
    n = read(fd[0], buf, sizeof buf);
    write(STDOUT_FILENO, buf, n);
    _exit(0);
}
close(fd[0]);                         /* chi scrive rinuncia al lato di lettura */
write(fd[1], "hello\n", 6);
close(fd[1]);                         /* ora, e solo ora, il figlio vede la fine del file */
\end{lstlisting}

\warning{La read che non restituisce mai zero}{
    Il genitore scrive una riga, chiude \texttt{fd[1]} e aspetta. Il figlio cicla sulla \texttt{read} e stampa. Non finisce
    mai. Perché? Perché il figlio \textbf{non ha chiuso il proprio} \texttt{fd[1]}. Un lato di scrittura è ancora aperto (il
    suo), quindi la pipe non arriva mai alla fine del file, la \texttt{read} si blocca, e il genitore aspetta un figlio che
    aspetta se stesso.
    \begin{itemize}
        \item La fine del file su una pipe significa \textbf{nessuno scrittore rimasto}, e il conto include la copia del
              lettore stesso.
        \item È la stessa regola del socket connesso: l'oggetto muore quando se ne va il suo ultimo descriptor, chiunque lo
              tenga.
        \item Il caso speculare è rumoroso invece che silenzioso: chi scrive in una pipe il cui lettore se n'è andato riceve
              \texttt{SIGPIPE}.
    \end{itemize}
}

\begin{table}[H]
\centering
\begin{tabular}{l l l}
\toprule
 & \textbf{L'altro lato è aperto} & \textbf{L'altro lato è chiuso} \\
\midrule
\texttt{read}, dati disponibili  & restituisce quel che c'è & restituisce quel che c'è \\
\texttt{read}, pipe vuota        & si blocca                & restituisce 0 \\
\texttt{write}, c'è posto        & copia e ritorna          & \texttt{SIGPIPE}, \texttt{EPIPE} \\
\texttt{write}, buffer pieno     & si blocca                & \texttt{SIGPIPE}, \texttt{EPIPE} \\
\bottomrule
\end{tabular}
\caption{Leggere e scrivere una pipe. Un lettore vede la fine solo quando i dati sono finiti \emph{e} non resta nessuno
         scrittore.}
\end{table}

\begin{itemize}
    \item Una \texttt{write} fino a \texttt{PIPE\_BUF} byte (4096 su Linux) è \textbf{atomica}: più scrittori non possono
          mescolarsi al suo interno.
    \item È la soglia dell'atomicità, non la dimensione della pipe: il buffer contiene 64 KB, e chi scrive si blocca quando
          è pieno.
\end{itemize}

\subsection{I figli che riferiscono al genitore}

Con le pipe la chat multiprocesso si può fare: ogni figlio passa al genitore le righe che legge, e il genitore le
ridistribuisce.

\begin{figure}[H]
\centering
\begin{tikzpicture}[every node/.style={font=\scriptsize\sffamily},
    ch/.style={draw=netblue!70!black, fill=netblue!8, rounded corners=2pt, align=center, minimum width=2.6cm,
               minimum height=0.7cm},
    a/.style={-{Stealth[length=1.8mm]}, thick},
    ax/.style={-{Stealth[length=1.8mm]}, thick, dashed, black!55}]
  \node[ch] (A) at (0,1.2)  {figlio A\\\texttt{fd 4} $\to$ client 1};
  \node[ch] (B) at (0,0)    {figlio B\\\texttt{fd 4} $\to$ client 2};
  \node[ch] (C) at (0,-1.2) {figlio C\\\texttt{fd 4} $\to$ client 3};
  \node[draw=netgreen!70!black, fill=netgreen!8, rounded corners=2pt, align=center, minimum height=1.0cm] (P) at (4.2,0)
    {\textbf{una pipe}\\molti scrittori\\un solo lettore};
  \node[draw=netred!70!black, fill=netred!6, rounded corners=2pt, align=center, minimum height=1.0cm] (G) at (8.0,0)
    {\textbf{genitore}\\legge, poi invia\\a tutti gli altri};
  \draw[a] (A.east) -- (P.west);
  \draw[a] (B.east) -- (P.west);
  \draw[a] (C.east) -- (P.west);
  \draw[a] (P.east) -- (G.west);
  \draw[ax] (G.south) -- (8.0,-2.3) -- (-1.9,-2.3) |- (A.west);
  \draw[ax] (-1.9,0) -- (B.west);
  \draw[ax] (-1.9,-1.2) -- (C.west);
  \node[text=black!60, fill=white] at (3.0,-2.3) {e indietro: un'altra pipe per ogni figlio};
\end{tikzpicture}
\caption{Ogni riga attraversa il kernel due volte ed è copiata una volta per destinatario, e il genitore deve stare in un
         ciclo tutto suo. Con i thread la stessa cosa è un \texttt{for} su un array.}
\end{figure}

Una pipe verso l'alto basta, perché le scritture sotto \texttt{PIPE\_BUF} non si mescolano; il ritorno verso il basso è la
metà costosa. E il genitore, che deve leggere dalla pipe e insieme accettare nuove connessioni, si trova già con il problema
del prossimo capitolo: aspettare su più descriptor alla volta.

\subsection{dup e dup2}

\begin{lstlisting}[language={[POSIX]C}, style=wnumbers, title={Due prototipi}]
int dup(int oldfd);
int dup2(int oldfd, int newfd);
\end{lstlisting}

\begin{itemize}
    \item Entrambe creano un \textbf{secondo descriptor sullo stesso oggetto}: stesso offset, stessa connessione, un unico
          contatore di riferimenti.
    \item \texttt{dup} prende il numero libero più basso; \texttt{dup2} prende il numero che indichiamo, chiudendolo prima se
          era aperto.
    \item Ed è proprio questo il punto: si può scegliere di atterrare su \textbf{0 o 1}.
\end{itemize}

\begin{figure}[H]
\centering
\begin{tikzpicture}[every node/.style={font=\scriptsize\sffamily},
    fd/.style={draw=black!70, fill=white, minimum width=1.5cm, minimum height=0.42cm, font=\scriptsize\ttfamily},
    fdc/.style={fd, draw=black!30, text=black!40, fill=black!4},
    ob/.style={draw=netblue!70!black, fill=netblue!8, rounded corners=2pt, minimum width=2.3cm, minimum height=0.6cm},
    a/.style={-{Stealth[length=1.6mm]}, thick}]
  \node[font=\small\bfseries\sffamily] at (0,1.15) {prima};
  \node[fd]  (p0) at (0,0.6)  {0};
  \node[fd]  (p1) at (0,0.15) {1};
  \node[fd]  (p2) at (0,-0.3) {2};
  \node[fdc] (p3) at (0,-0.75) {3 chiuso};
  \node[fd]  (p4) at (0,-1.2) {4};
  \node[ob] (T) at (3.6,0.3) {terminale};
  \node[ob, draw=netgreen!70!black, fill=netgreen!8] (K) at (3.6,-1.0) {connessione};
  \draw[a] (p0.east) -- (T.west);
  \draw[a] (p1.east) -- (T.west);
  \draw[a] (p2.east) -- (T.west);
  \draw[a] (p4.east) -- (K.west);
  \node[font=\small\bfseries\sffamily] at (7.2,1.15) {dopo \texttt{dup2} e \texttt{close(fd)}};
  \node[fd]  (d0) at (7.2,0.6)  {0};
  \node[fd]  (d1) at (7.2,0.15) {1};
  \node[fd]  (d2) at (7.2,-0.3) {2};
  \node[fdc] (d3) at (7.2,-0.75) {3 chiuso};
  \node[fdc] (d4) at (7.2,-1.2) {4 chiuso};
  \draw[a] (d0.west) -- (K.east);
  \draw[a] (d1.west) -- (K.east);
  \draw[a] (d2.west) -- (T.east);
\end{tikzpicture}
\caption{La tabella dei descriptor del figlio. Dopo i due \texttt{dup2} lo standard input e lo standard output sono la
         connessione; lo standard error resta sul terminale del server, dove finiscono i messaggi diagnostici.}
\end{figure}

\begin{lstlisting}[language={[POSIX]C}, style=wnumbers, title={Un programma che non sa di essere in rete}]
if (fork() == 0) {
    close(listen_fd);
    dup2(fd, STDIN_FILENO);               /* la connessione diventa lo standard input */
    dup2(fd, STDOUT_FILENO);              /* ... e lo standard output */
    if (fd > STDERR_FILENO) close(fd);    /* bastano i due duplicati */
    execl("/usr/local/bin/quoted", "quoted", (char *) NULL);
    perror("execl");                      /* si arriva qui solo se exec e' fallita */
    _exit(127);
}
close(fd);
\end{lstlisting}

Il programma dall'altra parte di \texttt{execl} legge con \texttt{scanf}, scrive con \texttt{printf} e non scopre mai da
dove arriva il suo input. È così che funziona un servizio avviato da \key{inetd}, ed è così che un server web esegue un
programma \key{CGI}.

\info{Il buffer di stdout dopo dup2}{
    Su un terminale \texttt{printf} svuota il buffer a ogni \texttt{\textbackslash n}; su un socket no, perché lo
    standard output non è più un terminale e la libreria C passa al buffer pieno. Il programma che risponde una riga e poi
    aspetta la successiva deve chiamare \texttt{fflush(stdout)}, oppure \texttt{setvbuf(stdout, NULL, \_IOLBF, 0)}
    all'inizio: altrimenti il client non riceve nulla finché il buffer non si riempie.
}

\subsection{Parentela, un nome, un indirizzo}

\begin{figure}[H]
\centering
\begin{tikzpicture}[every node/.style={font=\scriptsize\sffamily},
    pr/.style={draw=netblue!70!black, fill=netblue!8, rounded corners=2pt, align=center, minimum width=1.6cm,
               minimum height=0.6cm},
    md/.style={rounded corners=2pt, minimum width=2.6cm, minimum height=0.6cm, font=\scriptsize\ttfamily},
    ln/.style={thick, black!60}]
  \newcommand{\colonna}[7]{%
    \node[font=\small\bfseries\sffamily] at (#1,1.6) {#2};
    \node[pr] (l#1) at (#1-1.1,0.8) {#3};
    \node[pr] (r#1) at (#1+1.1,0.8) {#4};
    \node[md, draw=#6!70!black, fill=#6!8] (m#1) at (#1,-0.4) {#5};
    \draw[ln] (l#1.south) -- (m#1.north west);
    \draw[ln] (r#1.south) -- (m#1.north east);
    \node[align=center, text=black!70] at (#1,-1.4) {#7};}
  \colonna{0}{pipe}{genitore}{figlio}{buffer nel kernel}{netblue}
    {Devono essere parenti:\\il descriptor si eredita,\\non si può chiedere.}
  \colonna{5}{FIFO}{processo\\qualsiasi}{processo\\qualsiasi}{/tmp/upper}{netgreen}
    {Devono condividere un nome\\nel file system: una sola\\macchina, concordato prima.}
  \colonna{10}{socket}{host A}{host B}{143.225.5.3:5000}{netred}
    {Devono condividere un indirizzo\\e una porta, e nient'altro:\\nemmeno la macchina.}
\end{tikzpicture}
\caption{Tre modi in cui due processi si trovano, e a ogni passo devono avere meno cose in comune.}
\end{figure}

È la stessa forma cliente-servente tre volte: uno aspetta in un posto convenuto, l'altro ci va. Cambia solo che cosa
rende il posto raggiungibile.

\subsection{Una FIFO: server e client}

\dfn{FIFO}{
    Una \textbf{FIFO} (o \emph{named pipe}) è una pipe con un \textbf{nome nel file system}. Si crea con \texttt{mkfifo},
    si apre con \texttt{open} come un file qualsiasi, e da lì in poi si comporta esattamente come una pipe.
}

\begin{lstlisting}[language={[POSIX]C}, style=wnumbers, title={Un server FIFO e il suo client}]
/* server: crea il nome, poi aspetta su di esso */
mkfifo("/tmp/upper", 0666);                 /* se esiste gia' non e' un errore per noi */
int fd = open("/tmp/upper", O_RDONLY);      /* si blocca finche' uno scrittore non apre */
while ((n = read(fd, buf, sizeof buf - 1)) > 0) {
    buf[n] = '\0';
    for (int i = 0; i < n; i++) buf[i] = toupper((unsigned char) buf[i]);
    printf("%s\n", buf);
}

/* client: gli basta lo stesso nome */
int fd = open("/tmp/upper", O_WRONLY);      /* si blocca finche' un lettore non apre */
write(fd, argv[1], strlen(argv[1]));
close(fd);
\end{lstlisting}

\begin{lstlisting}[language=bash, title={Il nome nel file system}]
$ ./upper &
$ ls -l /tmp/upper
prw-r--r-- 1 utente utente 0 ott 19 10:02 /tmp/upper      <- p: una pipe con nome
$ ./shout ciao
CIAO
\end{lstlisting}

Due programmi, uno che aspetta e uno che chiede, e nessuna rete: il nome nel file system fa ciò che più avanti faranno un
indirizzo e una porta, e solo \texttt{unlink} lo toglie di mezzo.
\begin{itemize}
    \item Le due \texttt{open} si aspettano a vicenda: è il primo esempio di \textbf{appuntamento} fra due processi senza
          parentela.
    \item Al primo client che chiude, la \texttt{read} del server restituisce 0 e il server finisce. Per servirne altri lo
          si riapre in un ciclo, oppure lo si apre anche in scrittura (\texttt{O\_RDWR}): con uno scrittore sempre presente,
          la fine del file non arriva mai.
\end{itemize}

\subsection{Memoria condivisa attraverso una fork}

\begin{lstlisting}[language={[POSIX]C}, style=wnumbers, title={Un contatore che tutti i figli vedono}]
unsigned long *served = mmap(NULL, sizeof *served,
                             PROT_READ | PROT_WRITE,
                             MAP_SHARED | MAP_ANONYMOUS, -1, 0);
if (served == MAP_FAILED) { perror("mmap"); exit(EXIT_FAILURE); }
*served = 0;

/* ogni figlio creato da qui in poi vede questo stesso contatore */
\end{lstlisting}

\begin{itemize}
    \item \texttt{MAP\_ANONYMOUS} significa nessun file dietro: una pagina di memoria che i figli ereditano
          \textbf{condivisa} invece di copiarla.
    \item \texttt{MAP\_PRIVATE} al suo posto restituisce esattamente ciò che la \texttt{fork} dava già: una copia ciascuno.
    \item E \texttt{++*served} da due processi è la corsa critica che già conosciamo. Un mutex lì deve \textbf{vivere in
          questa pagina} ed essere creato con \texttt{PTHREAD\_PROCESS\_SHARED}.
\end{itemize}

\begin{lstlisting}[language={[POSIX]C}, style=wnumbers, title={Il mutex nella pagina condivisa}]
struct shared { pthread_mutex_t m; unsigned long served; };

struct shared *sh = mmap(NULL, sizeof *sh, PROT_READ | PROT_WRITE,
                         MAP_SHARED | MAP_ANONYMOUS, -1, 0);
pthread_mutexattr_t a;
pthread_mutexattr_init(&a);
pthread_mutexattr_setpshared(&a, PTHREAD_PROCESS_SHARED);
pthread_mutex_init(&sh->m, &a);           /* prima di ogni fork */
sh->served = 0;

/* in un figlio */
pthread_mutex_lock(&sh->m);
unsigned long me = ++sh->served;
pthread_mutex_unlock(&sh->m);
\end{lstlisting}

Un problema che con i thread non c'era: un figlio può \textbf{morire con il lock in mano}, e allora nessuno lo riaprirà
più. Un mutex creato anche con \texttt{PTHREAD\_MUTEX\_ROBUST} lo segnala al successivo che lo prende, con
\texttt{EOWNERDEAD}, invece di bloccarlo per sempre.

% ══════════════════════════════════════════════════════════════
\section{Un thread per client}
% ══════════════════════════════════════════════════════════════

Stesso server, stesse chiamate, \textbf{un solo spazio di indirizzi}: i descriptor stanno in un'unica tabella e lo stato
della chat è una variabile.

\begin{figure}[H]
\centering
\begin{tikzpicture}[every node/.style={font=\scriptsize\sffamily},
    th/.style={draw=netblue!70!black, fill=netblue!8, rounded corners=2pt, align=center, minimum width=2.5cm,
               minimum height=1.2cm}]
  \node[draw=black!60, thick, rounded corners, minimum width=10.2cm, minimum height=4.6cm] at (0,-0.35) {};
  \node[font=\small\bfseries\sffamily] at (0,1.55) {un processo, un solo spazio di indirizzi};
  \node[th] at (-3.4,0.5) {\textbf{thread 1}\\stack proprio\\\texttt{fd 4} client A};
  \node[th] at (0,0.5)    {\textbf{thread 2}\\stack proprio\\\texttt{fd 5} client B};
  \node[th] at (3.4,0.5)  {\textbf{thread 3}\\stack proprio\\\texttt{fd 6} client C};
  \node[draw=netgreen!70!black, fill=netgreen!10, rounded corners=2pt, align=center, minimum width=9.4cm,
        minimum height=0.9cm] at (0,-1.0)
    {\textbf{condiviso}: lo stato della chat, chi è connesso e che cosa inviargli\\una copia sola, raggiunta leggendo una
     variabile};
  \node[draw=netred!70!black, fill=netred!6, rounded corners=2pt, minimum width=9.4cm, minimum height=0.6cm] at (0,-2.1)
    {\textbf{thread principale}: \texttt{accept} su \texttt{fd 3}, un nuovo thread per ogni connessione accettata};
\end{tikzpicture}
\caption{Il server di chat con un thread per client. Ogni thread vede tutti i descriptor, quindi può scrivere a tutti i
         client: è ciò che mancava alla versione con i processi.}
\end{figure}

\begin{lstlisting}[language={[POSIX]C}, style=wnumbers, title={Il ciclo di accept che avvia un thread}]
int main(int argc, char *argv[]) {
    int listen_fd = listening_socket(argc > 1 ? argv[1] : "5000");
    signal(SIGPIPE, SIG_IGN);
    for (int i = 0; i < MAXC; i++) clients[i] = -1;

    for (;;) {
        int *fd = malloc(sizeof *fd);              /* l'argomento tutto suo del thread */
        if (fd == NULL) { sleep(1); continue; }
        *fd = accept(listen_fd, NULL, NULL);
        if (*fd < 0) { if (errno != EINTR) perror("accept"); free(fd); continue; }

        pthread_t tid;
        int err = pthread_create(&tid, NULL, worker, fd);
        if (err) { fprintf(stderr, "pthread_create: %s\n", strerror(err));
                   close(*fd); free(fd); continue; }
        pthread_detach(tid);                       /* nessuno lo aspettera' mai */
    }
}
\end{lstlisting}

\subsection{Lo stato che i thread condividono}

\begin{lstlisting}[language={[POSIX]C}, style=wnumbers, title={La lista dei client e il suo mutex}]
#define MAXC 64
static int clients[MAXC];                   /* socket connessi, -1 = posto libero */
static pthread_mutex_t roster = PTHREAD_MUTEX_INITIALIZER;

static int add_client(int fd) {
    int slot = -1;
    pthread_mutex_lock(&roster);
    for (int i = 0; i < MAXC; i++)
        if (clients[i] == -1) { clients[i] = fd; slot = i; break; }
    pthread_mutex_unlock(&roster);
    return slot;                            /* -1: la chat e' piena, si rifiuta */
}

static void del_client(int slot) {
    pthread_mutex_lock(&roster); clients[slot] = -1; pthread_mutex_unlock(&roster);
}
\end{lstlisting}

\begin{lstlisting}[language={[POSIX]C}, style=wnumbers, title={Il worker: registrarsi, servire, cancellarsi}]
static void *worker(void *arg) {
    int fd = *(int *) arg;
    free(arg);                              /* l'argomento era nostro: finito */

    int slot = add_client(fd);
    if (slot < 0) { write(fd, "chat is full\n", 13); close(fd); return NULL; }

    char buf[256];
    ssize_t n;
    while ((n = read(fd, buf, sizeof buf)) > 0)
        broadcast(slot, buf, n);            /* a tutti tranne me */

    del_client(slot);
    close(fd);
    return NULL;
}
\end{lstlisting}

\begin{lstlisting}[language={[POSIX]C}, style=wnumbers, title={Inviare a tutti gli altri}]
static void broadcast(int from, const char *msg, size_t len) {
    pthread_mutex_lock(&roster);
    for (int i = 0; i < MAXC; i++)
        if (clients[i] != -1 && i != from)
            write(clients[i], msg, len);    /* EPIPE: se ne accorgera' il thread di quel client */
    pthread_mutex_unlock(&roster);
}
\end{lstlisting}

\begin{itemize}
    \item Questa è \textbf{tutta la differenza} dalla versione multiprocesso: un \texttt{for} su una variabile che ogni
          thread può vedere.
    \item Il mutex non è facoltativo: un altro thread potrebbe aggiungersi a \texttt{clients} mentre questo ciclo la legge.
    \item Ed è tenuto \textbf{durante una write}, che è proprio l'unica cosa che avevamo detto di non fare mai.
\end{itemize}

\warning{Il client lento che congela la chat}{
    Un client smette di leggere: la sua finestra si riempie, e si riempie anche il buffer del suo socket. Che cosa succede
    agli altri diciannove? La \texttt{write} verso quel client \textbf{si blocca}, con il mutex \texttt{roster} in mano.
    Ogni altro thread si ferma al suo \texttt{broadcast} successivo: un solo client silenzioso ha fermato l'intera chat.
}

Rilasciare il lock prima di scrivere non lo risolve: fra la copia della lista e la \texttt{write} quel posto può essere
chiuso, e il suo numero di descriptor riusato da un nuovo client.

\begin{table}[H]
\centering
\begin{tabular}{c l l l}
\toprule
\textbf{Passo} & \textbf{Thread A (broadcast)} & \textbf{Thread B} & \textbf{Thread principale} \\
\midrule
1 & copia la lista: posto 3 = \texttt{fd 7} & & \\
  & rilascia il lock & & \\
2 & & il client 3 se ne va: & \\
  & & \texttt{del\_client(3)}, \texttt{close(7)} & \\
3 & & & \texttt{accept} restituisce di nuovo 7 \\
4 & \texttt{write(7, ...)}: al nuovo arrivato & & \\
\bottomrule
\end{tabular}
\caption{Perché non basta copiare la lista e scrivere senza lock. Il kernel riassegna sempre il descriptor libero più
         basso, quindi il riuso non è un caso raro: la riga va alla connessione sbagliata, di rado e mai mentre si guarda.}
\end{table}

Nessuna delle due strade è la risposta. Un server vero \textbf{non si blocca mai in un worker}: ogni client ha una
\key{coda in uscita}, e qualcuno la svuota quando il socket è pronto. È l'argomento del prossimo capitolo.

% ══════════════════════════════════════════════════════════════
\section{Processi o thread, per un server}
% ══════════════════════════════════════════════════════════════

\begin{table}[H]
\centering
\begin{tabular}{l l l}
\toprule
 & \textbf{Un processo per client} & \textbf{Un thread per client} \\
\midrule
stato condiviso             & deve portarlo il kernel          & una variabile \\
il crash di un client       & uccide il suo figlio             & uccide il server \\
un descriptor dimenticato   & se ne va con il figlio           & resta per sempre \\
costo di avvio              & uno spazio di indirizzi          & uno stack \\
\texttt{close} su un socket & chiude solo quella copia         & lo chiude per tutti \\
debugging                   & un flusso, molti processi        & un processo, molti flussi \\
\bottomrule
\end{tabular}
\caption{I processi pagano per essere isolati; i thread pagano per condividere. Una chat vuole condividere, un server che
         esegue codice non fidato vuole l'isolamento.}
\end{table}

\Needspace{16\baselineskip}
\warning{Un worker dereferenzia un puntatore nullo}{
    Un bug nel codice che serve un client: quanto costa? \textbf{Con i processi}: muore la connessione di quel client, il
    genitore raccoglie un figlio terminato da un segnale, e diciannove conversazioni continuano. \textbf{Con i thread}: il
    processo riceve \texttt{SIGSEGV}, e ogni connessione muore con lui.
    \begin{itemize}
        \item La stessa asimmetria vale per un buffer perso, uno heap corrotto, un descriptor che nessuno chiude: con i
              processi tutto finisce a \texttt{\_exit}.
        \item È per questo che la scelta si fa ancora in entrambi i modi: Apache offre un modello a processi e uno a thread,
              e quello a processi si sceglie quando il codice non è fidato.
    \end{itemize}
}

\subsection{Quanto costa «uno per client»}

\begin{itemize}
    \item Entrambe le versioni pagano \textbf{per connessione}, al momento della connessione: uno spazio di indirizzi o uno
          stack, creati e distrutti per un client che magari dice una parola.
    \item Entrambe mantengono quel costo finché il client resta, che parli o no; e un client di chat tace quasi sempre.
    \item Diecimila connessioni inattive significano diecimila stack e diecimila oggetti del kernel, per non fare nulla.
    \item E una raffica di connessioni brevi passa più tempo a creare e distruggere che a servire.
\end{itemize}

Entrambi i limiti indicano la stessa direzione: smettere di creare \textbf{un flusso di controllo per client}. Ne serve uno
che osserva molte connessioni e si accorge di quali sono pronte, più un insieme fisso di worker tenuti occupati da una
coda.

% ══════════════════════════════════════════════════════════════
\section{Esercizi}
% ══════════════════════════════════════════════════════════════

\begin{esercizio}{Una pipeline fatta a mano}
Scrivere un programma che esegua l'equivalente di \texttt{ls | wc -l} usando soltanto \texttt{pipe}, \texttt{fork},
\texttt{dup2}, \texttt{execlp} e \texttt{wait}. Poi togliere dal genitore la \texttt{close(fd[1])}: che cosa succede, e
perché?

\svolgimento
Due figli, uno per programma: il primo mette il lato di scrittura sul proprio standard output, il secondo il lato di
lettura sul proprio standard input. Poi ciascuno chiude \textbf{entrambi} gli originali, perché i duplicati bastano.

\begin{lstlisting}[language={[POSIX]C}, style=wnumbers, title={pipeline.c}]
int fd[2];
if (pipe(fd) < 0) { perror("pipe"); exit(EXIT_FAILURE); }

if (fork() == 0) {                      /* primo figlio: ls */
    dup2(fd[1], STDOUT_FILENO);
    close(fd[0]); close(fd[1]);
    execlp("ls", "ls", (char *) NULL);
    perror("execlp"); _exit(127);
}
if (fork() == 0) {                      /* secondo figlio: wc -l */
    dup2(fd[0], STDIN_FILENO);
    close(fd[0]); close(fd[1]);
    execlp("wc", "wc", "-l", (char *) NULL);
    perror("execlp"); _exit(127);
}
close(fd[0]); close(fd[1]);             /* il genitore non usa la pipe */
while (wait(NULL) > 0)
    ;
\end{lstlisting}

Senza la \texttt{close(fd[1])} nel genitore, \texttt{ls} stampa, esce, e \texttt{wc} \textbf{non termina mai}: un lato di
scrittura è ancora aperto (quello del genitore), quindi la sua \texttt{read} non restituisce mai 0. E il genitore, fermo
in \texttt{wait}, aspetta un figlio che aspetta lui. È il caso del riquadro sulla \texttt{read} che non restituisce mai
zero, con tre processi invece di due. La shell fa esattamente queste chiamate per ogni \texttt{|}.
\end{esercizio}

\begin{esercizio}{Iterazione 2: una chat multiprocesso}
Si riparte dal server sequenziale: corretto, e inutilizzabile appena si connette un secondo client, quello lasciato ad
aspettare in \texttt{ss -tn}.
\begin{itemize}
    \item \textbf{Che cosa cambia}: solo il ciclo di \texttt{accept}. La \texttt{fork} dopo \texttt{accept}, le due
          \texttt{close}, \texttt{\_exit} nel figlio, e un handler di \texttt{SIGCHLD} installato prima del primo client.
    \item \textbf{Che cosa no}: \texttt{serve()} resta intatta. Se ci si ritrova a modificarla, si sta risolvendo il
          problema sbagliato.
\end{itemize}
Connettere venti client e poi ucciderne uno. Dimostrare entrambe le cose: venti figli, e nessuno zombie. E che cosa dice
il contatore al ventesimo client?

\svolgimento
Il \texttt{main} è quello del capitolo: \texttt{reap} installato con \texttt{sigaction}, \texttt{SIGPIPE} ignorato, e il
ciclo con la \texttt{fork}. La prova si fa da un secondo terminale:

\begin{lstlisting}[language=bash, title={Venti figli, nessuno zombie}]
$ ./chatd 5000 &
$ for i in $(seq 1 20); do sleep 60 | nc 127.0.0.1 5000 & done
$ ps --ppid $(pgrep -o chatd) --no-headers | wc -l
20
$ kill %2                                   <- termina uno dei client
$ ps --ppid $(pgrep -o chatd) --no-headers | wc -l
19
$ ps --ppid $(pgrep -o chatd) -o stat --no-headers | grep -c Z
0
\end{lstlisting}

Quando il client muore, la \texttt{read} del suo figlio restituisce 0, il figlio esce, e il \texttt{SIGCHLD} lo fa
raccogliere subito dal genitore: i figli scendono a diciannove e nessuno resta \texttt{Z}. Togliendo la
\texttt{sigaction}, lo stesso comando mostrerebbe un \texttt{Z} per ogni client che se n'è andato.

Il ventesimo client si sente rispondere \texttt{1:} alla sua prima riga, come tutti gli altri. Nell'iterazione 1 il
contatore era locale a \texttt{serve}; ma anche spostandolo in una globale non cambierebbe nulla, perché ogni figlio ne
riceve una copia al momento della \texttt{fork}. Per un contatore comune servirebbe la pagina \texttt{mmap} con il suo
mutex.
\end{esercizio}

\begin{esercizio}{Iterazione 3: la stessa chat con i thread}
Si riparte dall'iterazione 2, che serve tutti e non avvisa nessuno: una riga scritta da un client arriva solo al figlio che
l'ha letta.
\begin{itemize}
    \item \textbf{Che cosa cambia}: \texttt{pthread\_create} e \texttt{pthread\_detach} nel ciclo, una \texttt{malloc} per
          connessione per l'argomento, e il worker del capitolo.
    \item \textbf{Che cosa è nuovo}: la lista dei client e il suo mutex, perché una riga di un client arrivi a tutti gli
          altri. È la prima versione che è davvero una chat.
\end{itemize}
Poi misurare la memoria per client connesso, con 50 client, in entrambe le versioni. Quale costa di più per client, e quale
si sceglierebbe comunque se \texttt{serve()} eseguisse codice scritto da altri?

\svolgimento
Il worker del capitolo inoltra quello che \texttt{read} restituisce, che non è ancora una riga. Si recupera il buffer di
riga dell'iterazione 1, e \texttt{broadcast} parte solo a riga completa:

\begin{lstlisting}[language={[POSIX]C}, style=wnumbers, title={chatd.c: il worker inoltra righe intere}]
static void *worker(void *arg) {
    int fd = *(int *) arg;
    free(arg);
    int slot = add_client(fd);
    if (slot < 0) { write(fd, "chat is full\n", 13); close(fd); return NULL; }

    char buf[512], line[512];
    size_t len = 0;
    ssize_t n;
    while ((n = read(fd, buf, sizeof buf)) > 0)
        for (ssize_t i = 0; i < n; i++) {
            if (len < sizeof line - 1) line[len++] = buf[i];
            if (buf[i] == '\n') { broadcast(slot, line, len); len = 0; }
        }
    del_client(slot);
    close(fd);
    return NULL;
}
\end{lstlisting}

\begin{lstlisting}[language=bash, title={La misura}]
$ for i in $(seq 1 50); do sleep 300 | nc 127.0.0.1 5000 & done
$ ps -o pid,nlwp,rss,cmd -p $(pgrep -o chatd)           <- iterazione 3: un processo, 51 thread
$ ps -o rss= --ppid $(pgrep -o chatd) | paste -sd+ | bc  <- iterazione 2: somma dei 50 figli
\end{lstlisting}

Con i thread la memoria residente cresce di poco per client: ogni stack riserva 8 MB di spazio virtuale, ma diventano
memoria vera solo le poche pagine che il thread tocca. Con i processi ogni figlio ha la propria tabella delle pagine, le
proprie strutture nel kernel e le proprie pagine di stack e di dati; il codice e le librerie restano condivise, ed è per
questo che \textbf{sommare l'RSS dei figli sovrastima}: conta le pagine condivise una volta per figlio. La misura onesta è
il PSS, in \texttt{/proc/PID/smaps\_rollup}, che divide ogni pagina condivisa fra chi la usa. Anche così, il processo costa
più del thread.

Ma se \texttt{serve()} eseguisse codice di altri si sceglierebbero comunque i \textbf{processi}: un bug in quel codice
ucciderebbe un figlio solo, invece dell'intero server.
\end{esercizio}

\Needspace{20\baselineskip}
% ══════════════════════════════════════════════════════════════
\section{Riepilogo}
% ══════════════════════════════════════════════════════════════

\riepilogo{
\begin{itemize}
    \item Un \textbf{server concorrente} è lo stesso servizio in un ciclo diverso: \texttt{fork} o \texttt{pthread\_create}
          subito dopo \texttt{accept}, e il genitore o il thread principale di nuovo in \texttt{accept}.
    \item Con i processi il prezzo si paga subito: i figli non si vedono, quindi ogni cosa condivisa attraversa il kernel
          (una pipe, o una pagina di \texttt{mmap} con un lock dentro).
    \item Una \textbf{pipe} unisce parenti, una \textbf{FIFO} chiunque abbia il nome, un \textbf{socket} chiunque abbia
          l'indirizzo: la stessa forma cliente-servente tre volte.
    \item \texttt{dup2} mette una connessione sui descriptor 0 e 1, e un programma che non sa nulla di reti diventa un
          servizio di rete.
    \item Con i thread la condivisione è gratuita e la sicurezza no: la lista dei client vuole un mutex, e tenerlo durante
          una \texttt{write} consegna l'intera chat al client più lento.
    \item Isolamento contro condivisione, e un costo che nessuno dei due evita: \textbf{un flusso di controllo per
          connessione}.
\end{itemize}
}
